\documentclass[11pt]{article}
\usepackage[margin=1in]{geometry}
\usepackage[T1]{fontenc}
\usepackage{lmodern,amsmath,amssymb,amsthm,mathtools,booktabs,array,longtable,graphicx}
\usepackage{microtype,url,hyperref,placeins}
\hypersetup{hidelinks,pdftitle={Contextuality and Exact Information Resources over Finite Chain Rings}}
\newtheorem{theorem}{Theorem}[section]
\newtheorem{lemma}[theorem]{Lemma}
\newtheorem{proposition}[theorem]{Proposition}
\newtheorem{corollary}[theorem]{Corollary}
\theoremstyle{definition}\newtheorem{definition}[theorem]{Definition}
\theoremstyle{remark}\newtheorem{remark}[theorem]{Remark}
\newcommand{\F}{\mathbb F}
\newcommand{\A}{\F_2[u]/(u^4)}
\DeclareMathOperator{\GL}{GL}
\DeclareMathOperator{\Mat}{Mat}
\DeclareMathOperator{\rank}{rank}
\DeclareMathOperator{\diag}{diag}
\DeclareMathOperator{\Span}{span}
\setlength{\emergencystretch}{3em}
\title{Contextuality and Exact Information Resources\\over Finite Chain Rings}
\author{Brian Theory\\\small B-Theory}
\date{26 September 2026}
\begin{document}
\maketitle
\begin{abstract}
We study nonzero bipartite matrices over a finite commutative chain ring using all primitive projective local bases, setting-indexed outcomes and nonzero-amplitude support. Two Smith invariants determine the support type: the number $r$ of nonzero factors and the multiplicity $h$ of the least exponent. The model is local when $r=1$, logically but not strongly contextual when $r\geq2$ and $h=1$, and strongly contextual when $h\geq2$. A residue-hyperplane criterion and a valuation-uniform Hardy construction give the classification. For a square $d$-dimensional resource, invertible local orbit encoding and one arbitrary complete joint basis readout admit exactly $dh$ messages. Universal exact raw-vector teleportation, with a complete joint covector basis and unrestricted invertible corrections, instead requires an invertible resource. We prove the existence of complete invertible-matrix analyzers over finite local rings and exhibit nonprimitive resources separating maximal orbit coding from teleportation. For primitive resources, maximal coding and universal exact raw teleportation coincide. Tensor choice and preparation closure delimit the interpretation. Finite checks accompany the proofs; no physical probability model or priority claim follows from them.
\end{abstract}

\section{Introduction}
Replacing complex amplitudes by a finite scalar system makes it possible to separate algebraic support from probability. In modal quantum theory, an effect is possible when its contraction with a state is nonzero. Field models already support entanglement, nonlocal support and communication protocols \cite{Schumacher2010,Schumacher2012}. Passing from a field to a ring introduces a different question: which aspects of a resource survive when nonzero amplitudes can annihilate one another?

This paper answers that question for static bipartite resources over finite commutative chain rings, with complete local basis families. The classification uses two elementary pieces of Smith data. At least two nonzero Smith factors are necessary and sufficient for a Hardy obstruction. Leading multiplicity detects whether any global outcome assignment survives. Under a separately stated exact communication task, the same leading multiplicity measures the number of linearly distinguishable orbit messages. Universal exact teleportation distinguishes yet another property, invertibility.

The contribution is the complete chain-ring classification and the comparison of contextuality, exact orbit coding and raw-vector teleportation through related Smith invariants under explicit contracts. Smith theory, the language of support contextuality, Hardy implications and modal protocols are established ingredients. The results do not propose a physical theory, and do not infer a Born rule. In particular, allowing all nonzero ring vectors is useful for static resource questions but does not define a preparation theory closed under arbitrary balanced tensor products.

The proofs are self-contained. The source comparison below identifies established ingredients and the hypotheses that distinguish the present tasks. It makes no claim of historical firstness. A standalone computational supplement supplies finite checks with their exact scope; an optional Boolean-interface note is separate from the article and is not a dependency of its results.

\subsection{Relation to established work}
Schumacher and Westmoreland's two-mobit model already gives a nonlocal support table and a four-message orbit code \cite{Schumacher2010}; their later treatment develops generalized field effects and modal probability-completion questions \cite{Schumacher2012}. These supply background, not new claims here. The support hierarchy follows the global-section formulation of Abramsky and Brandenburger \cite{AbramskyBrandenburger2011}. Our settings are local bases. We do not additionally require the same effect ray to carry a context-independent truth value across distinct bases.

That distinction matters for comparisons with modal noncontextuality and free-will arguments \cite{SchumacherFreeWill} and the uncolorability of integer-matrix idempotents \cite{BenZvi2017}. The latter concerns compatible idempotent decompositions and transports to arbitrary scalar rings; it is not a Smith classification of a fixed bipartite state with basis-indexed outcomes. Finite-chain-ring module structure and projective ring geometry are likewise classical \cite{HonoldLandjev2000,Byrne2021}. The use of primitive rows is a standard local-ring construction. Our assertions concern their complete bilinear support tables and two specified information tasks.

State admissibility is a substantive difference from other ring models. In de Beaudrap's generic state space, the coordinates must generate the unit ideal \cite{deBeaudrap2014}. Over a local ring this means that at least one coordinate is a unit. Our nonprimitive nilpotent separator is therefore excluded from that state space; its conclusions belong to the static all-nonzero-vector contract used here. Categorical finite-field protocols can also use a different measurement notion: Varughese's Example~8.3 uses three Bell outcomes on a nine-dimensional pair \cite{Varughese2009}. It does not provide the complete nine-row covector analyzer required by our teleportation contract.



The matrix-spanning ingredient is classical: nonsingular-summand results go back to Wolfson and Zelinsky, and Henriksen proves the stronger statement that every square matrix of size greater than one over a unital ring is a sum of three units \cite{Zelinsky1954,Henriksen1974}. Lemma~\ref{lem:span} supplies a short proof sufficient for our use, followed by residue-basis lifting. Neither this spanning fact nor the existence of modal protocols is asserted as a new general result.

Finite-chain-ring coding also uses module invariants, but its optimization problem must be distinguished. Feng et al.\ study a stochastic channel $Y=AX+BE$ with random transfer and error matrices, deriving capacity results under specified distributions \cite{Feng2014}. Here a sender controls $U$ in $UM$ for a fixed resource $M$, and a single complete joint basis readout must identify the message at every supported outcome. A direct consequence of the former channel definition makes the distinction concrete: if $A$ is uniform in $\GL_d(K)$ and noise is absent, then $AUM$ has the same distribution for every $U\in\GL_d(K)$. This observation is our comparison of the two tasks, not a claim made in that source about the present theorem.

\begin{table}[ht]\centering\small
\begin{tabular}{>{\raggedright\arraybackslash}p{.19\linewidth}>{\raggedright\arraybackslash}p{.25\linewidth}>{\raggedright\arraybackslash}p{.43\linewidth}}
\toprule
Source and domain & Admissible objects & Operations, readout and conclusion\\\midrule
Modal theory \cite{Schumacher2010,Schumacher2012}; fields & Nonzero state vectors; composite vector spaces & Reversible linear maps and basis effects; nonzero contraction gives possibility. Includes finite-field coding and teleportation examples.\\[3pt]
Ring computation \cite{deBeaudrap2014}; commutative rings & Generic vectors whose coordinates generate the unit ideal & Model-dependent algebraic transformations and outcome rules; this admissibility excludes our nonprimitive separator.\\[3pt]
Chain-ring codes \cite{HonoldLandjev2000,Byrne2021} & Submodules and their module types & Code equivalence, geometry and enumeration; no fixed-resource joint support readout or raw teleportation task.\\[3pt]
Matrix channels \cite{Feng2014}; chain rings & Input matrices in a prescribed ambient module & Random $A,B,E$ in $Y=AX+BE$; stochastic decoding and capacity, rather than a controlled orbit with disjoint nonzero supports.\\[3pt]
This article; chain rings & Every nonzero static matrix, including nonprimitive ones & All local bases for support; full $\GL_d(K)$ encoders and one complete joint decoder for exact coding; complete covector analysis and invertible corrections for exact raw recovery. Conclusions: Smith trichotomy, $dh$ messages, and invertibility criterion.\\
\bottomrule
\end{tabular}
\caption{Comparison by scalar domain, admissibility and task. The support hierarchy itself is the established global-section hierarchy \cite{AbramskyBrandenburger2011}. Theorem~\ref{thm:teleport} also applies to finite commutative local rings that are not chain rings.}\label{tab:prior}
\end{table}

\FloatBarrier
\section{Scalar, state and measurement contracts}\label{sec:contract}
Let $K$ be a finite commutative chain ring with identity, maximal ideal $J=(\pi)$, residue field $k=K/J$ of size $q$, and nilpotency index $s$, so $\pi^s=0$ and $\pi^{s-1}\ne0$. The field case has $s=1$. Bars denote reduction modulo $J$. An element is a unit exactly when its residue is nonzero. A square matrix is invertible exactly when its residue is invertible, by the determinant criterion.

Every nonzero element is a unit times a unique power $\pi^a$, $0\leq a<s$. Repeatedly moving an entry of minimal valuation to the first position and using elementary row and column operations gives Smith form
\[
 PMQ^T=\diag(\pi^{a_1},\ldots,\pi^{a_r},0,\ldots),\qquad
 0\leq a_1\leq\cdots\leq a_r<s,
\]
with $P\in\GL_m(K)$ and $Q\in\GL_n(K)$. Clearing a pivot is possible because it divides all remaining entries; induction gives the diagonal form. The invariant exponents can equivalently be recovered from the module type \cite{HonoldLandjev2000,Byrne2021}. We write $r$ for the number of nonzero factors and
\[
 a=a_1,\qquad h=\#\{i:a_i=a\}.
\]
Writing $M=\pi^a N$ involves a choice of preimage, not division by a unit. Its residue $\bar N$ is well-defined: entrywise, the ambiguity lies in $\operatorname{Ann}(\pi^a)=(\pi^{s-a})\subseteq J$ (and is zero when $a=0$). Thus $h=\rank_k\bar N$. In contrast, $\bar M$ is zero whenever $a>0$; its rank is then not the leading multiplicity.

\begin{definition}[Complete-basis support]
A primitive row in $K^m$ has a unit coordinate. Its projective ray is its orbit under unit scaling. A local setting is an unordered basis of primitive projective rows, represented by an invertible matrix. For a nonzero resource $M\in K^m\otimes_K K^n\cong\Mat_{m\times n}(K)$, the outcome pair represented by rows $x,y$ is supported exactly when
\[
 xMy^T\ne0.
\]
All such local settings are admitted, and local reversible changes are the full groups $\GL_m(K)$ and $\GL_n(K)$.
\end{definition}

Unit rescaling does not change support. Each context has a supported event because invertible changes cannot send nonzero $M$ to zero. The support marginal for $x$ is independent of Bob's basis: some coordinate of $xMB^T$ is nonzero if and only if $xM\ne0$. Hence these support tables satisfy possibilistic no-signalling, without any assigned real probabilities.

A \emph{global assignment} chooses one outcome for each local setting, with every chosen cross pair supported. A model is \emph{local} here if every supported event extends to such an assignment. It is \emph{logically contextual} if some event does not extend, and \emph{strongly contextual} if no global assignment exists. This is a support-level definition; ``local'' is not a statement about a particular probability distribution.

Two further tasks use the same scalars but different measurements. Their distinctions are essential.
\begin{center}\small
\begin{tabular}{p{.15\linewidth}p{.74\linewidth}}
\toprule
Task & Contract \\\midrule
Support & Nonzero static $M$; balanced tensor over $K$; all primitive projective local bases; nonzero contraction means possible. \\
Orbit coding & Square $M$, $d\geq2$; encodings $UM$ with $U\in\GL_d(K)$; one invertible $K$-linear joint decoder on $d^2$ coordinates; every possible outcome identifies the message. No discard and reprepare. \\
Teleportation & Every nonzero raw input, including correlations with an untouched reference; complete basis of primitive joint covectors; all possible branches corrected exactly by allowed invertible maps. No selected successful branch or postselection. \\\bottomrule
\end{tabular}
\end{center}
For support and coding, multiplying a whole resource by a unit is immaterial. Teleportation below demands actual vector recovery and does not silently replace this by an unspecified projective task.



\subsection{Resource map and a worked comparison}
Table~\ref{tab:map} previews the consequences of Theorems~\ref{thm:trichotomy}, \ref{thm:code} and~\ref{thm:teleport} for square resources. Primitivity of $M$, viewed as a vector of entries, is equivalent to $a=0$.
\begin{table}[ht]\centering\small
\begin{tabular}{ccl}
\toprule
Smith region & Orbit messages & Support and raw teleportation\\\midrule
$r=1$ & $d$ & Local; singular for $d\geq2$\\
$r\geq2,\ h=1$ & $d$ & Logical, not strong; singular\\
$2\leq h<d$ & $dh$ & Strong; singular\\
$h=d,\ a>0$ & $d^2$ & Strong; nonprimitive, singular\\
$h=d,\ a=0$ & $d^2$ & Strong; invertible, universally teleporting\\\bottomrule
\end{tabular}
\caption{Square resources under the stated contracts, $d\geq2$. Zero is excluded.}\label{tab:map}
\end{table}

For $A=\A$, compare the four resources in Table~\ref{tab:examples}. In $M=\diag(u,u)=uI_2$, a choice of leading preimage is $N=I_2$, so $a=1$ and $h=2$ even though $\bar M=0$. All four orbit messages remain distinguishable, but the nonzero input $(u^3,0)^T$ is annihilated by $M$ and cannot be universally recovered. By contrast, $\diag(1,u^2)$ is primitive and has two nonzero factors, but only one leading factor; it has a Hardy obstruction and also a global assignment. These examples distinguish $r$, $h$ and invertibility without changing the scalar ring.
\begin{table}[ht]\centering\small
\begin{tabular}{lccclcc}
\toprule
$M$ over $A$ & $r$ & $a$ & $h$ & Support & Messages & Teleports\\\midrule
$\diag(1,0)$ & 1 & 0 & 1 & Local & 2 & No\\
$\diag(1,u^2)$ & 2 & 0 & 1 & Logical, not strong & 2 & No\\
$\diag(u,u)$ & 2 & 1 & 2 & Strong & 4 & No\\
$I_2$ & 2 & 0 & 2 & Strong & 4 & Yes\\\bottomrule
\end{tabular}
\caption{A worked comparison over one chain ring. Teleportation means universal exact raw-vector recovery.}\label{tab:examples}
\end{table}

\section{Complete-basis contextuality}\label{sec:classification}
The first lemma requires only that $K$ be finite, commutative and local.
\begin{lemma}[Residue-hyperplane criterion]\label{lem:hyperplane}
A complete-basis support model has a global assignment if and only if there are hyperplanes $H_A\subset k^m$ and $H_B\subset k^n$ such that $xMy^T\ne0$ for every pair of primitive lifts whose residues lie outside $H_A$ and $H_B$.
\end{lemma}
\begin{proof}
Let $U_A$ be the union of the rays selected by a global assignment. Its complement contains no basis, since the corresponding setting would have to select a ray in that complement. If the reductions of the complement spanned $k^m$, one could select a residue basis from them; its lifts would form an invertible basis over $K$. Thus those reductions lie in a proper subspace, which is contained in a hyperplane $H_A$. Every primitive lift of a residue vector outside $H_A$ belongs to $U_A$. Construct $H_B$ in the same way. Every pair from $U_A\times U_B$ is selected in some pair of settings, and is therefore supported.

Conversely, any basis has a row reducing outside a given hyperplane. Select one in every setting on each side. The stated condition makes every chosen cross pair supported, giving a global assignment.
\end{proof}

\begin{lemma}[Uniform Hardy witness]\label{lem:hardy}
For $D=\diag(p,q_0)$ with $p=\pi^a$, $q_0=\pi^b\ne0$ and $a\leq b$, put $c=\pi^{b-a}$. The local bases
\[
 F=\begin{pmatrix}0&1\\1&0\end{pmatrix},\quad
 X=\begin{pmatrix}1&0\\1&1\end{pmatrix},\quad
 C=\begin{pmatrix}1&0\\-c&1\end{pmatrix}
\]
give a supported event with no global extension, using $F,C$ for Alice and $F,X$ for Bob.
\end{lemma}
\begin{proof}
All three determinants are units, and $pc=q_0$. Direct multiplication gives
\begin{align*}
 FDF^T&=\begin{pmatrix}q_0&0\\0&p\end{pmatrix},&
 FDX^T&=\begin{pmatrix}0&q_0\\p&p\end{pmatrix},\\
 CDX^T&=\begin{pmatrix}p&p\\-q_0&0\end{pmatrix},&
 CDF^T&=\begin{pmatrix}0&p\\q_0&-q_0\end{pmatrix}.
\end{align*}
Index outcomes by $0,1$. The forcing chain from the supported seed is
\[
(F_A=0,F_B=0)\ \Longrightarrow\ X_B=1
\ \Longrightarrow\ C_A=0\ \Longrightarrow\ F_B=1,
\]
a contradiction. Each implication excludes an outcome using, successively, the zero entries in $FDX^T$, $CDX^T$ and $CDF^T$. The support strings, in row order, are $1001,0111,1110,0111$. The minus sign cannot be dropped outside characteristic two.
\end{proof}

\begin{theorem}[Smith trichotomy]\label{thm:trichotomy}
For nonzero $M\in\Mat_{m\times n}(K)$, the complete-basis support type is
\[
\begin{array}{rcl}
r=1&\Longleftrightarrow&\text{local},\\
r\geq2,\ h=1&\Longleftrightarrow&\text{logical but not strong},\\
h\geq2&\Longleftrightarrow&\text{strong}.
\end{array}
\]
\end{theorem}
\begin{proof}
Invertible left and right changes merely relabel the complete basis families, so take Smith form.

Suppose $r=1$, with sole nonzero entry $p$ in position $(1,1)$. Retain the outcomes of any supported seed event $(x,y)$. In every other setting choose a row with unit first coordinate. A new--new pair has contraction equal to $p$ times a unit. The seed condition $px_1y_1\ne0$ implies $px_1\ne0$ and $py_1\ne0$, so a seed--new pair is also supported. The seed therefore extends.

If $h=1$, choose a row with unit first coordinate in every setting. In the divided matrix, each selected contraction is a product of units plus terms in the radical, hence is a unit. Multiplication by the nonzero $\pi^a$ gives a supported cross pair. A global assignment exists.

If $h\geq2$, write $M=\pi^a M_0$, so $\bar M_0$ has rank $h$. Suppose the hyperplanes of Lemma~\ref{lem:hyperplane} existed, and write $H_B=\ker\beta$. The vectors outside $H_A$ span $k^m$. Since the row space of $\bar M_0$ has dimension at least two, some $x\notin H_A$ has $w=x\bar M_0$ outside the line spanned by $\beta$. There is $y\in\ker w$ with $\beta(y)\ne0$. Lift $x,y$ to rows $\tilde x,\tilde y$. The row $\tilde xM_0$ has a unit coordinate, whereas its contraction with $\tilde y$ lies in $J$. Subtracting that contraction divided by the unit from the corresponding coordinate of $\tilde y$ cancels it exactly and changes no residue. This yields lifts outside both hyperplanes with zero contraction, a contradiction.

Finally, if $r\geq2$, embed Lemma~\ref{lem:hardy} in the first two nonzero Smith coordinates and extend its bases with standard rows. The embedded matrices are block diagonal. At each forcing step, the already selected row is supported only in the first two coordinates and therefore has zero contraction with every added coordinate outcome on the other side. No additional outcome can escape the forcing chain. Thus the same seed cannot extend, even in higher dimensions. These observations exhaust all nonzero cases.
\end{proof}

There are $q^{sm}-q^{(s-1)m}$ primitive rows in $K^m$. Unit scaling acts freely, and $|K^\times|=q^{s-1}(q-1)$. Consequently the numbers of rays and unordered projective bases are
\[
 q^{(s-1)(m-1)}\frac{q^m-1}{q-1},\qquad
 \frac{q^{m^2(s-1)}|\GL_m(k)|}{|K^\times|^m m!}.
\]
For $K=\A$ and $m=2$, these are $24$ rays and $192$ bases. These counting facts are background; enumeration is unnecessary for the classification proof.

\section{Exact orbit coding and teleportation}\label{sec:tasks}
The following classical spanning fact is included with a constructive proof; stronger units-generation results are available in \cite{Zelinsky1954,Henriksen1974}.
\begin{lemma}[Invertible matrices span]\label{lem:span}
For a field $k$ and $d\geq2$, $\GL_d(k)$ spans $\Mat_d(k)$. Over a finite commutative local ring $K$, there is a $K$-basis of $\Mat_d(K)$ consisting entirely of invertible matrices.
\end{lemma}
\begin{proof}
For $i\ne j$, $E_{ij}=(I+E_{ij})-I$. For a diagonal unit $E_{ii}$, choose a permutation matrix $P$ with $Pe_j=e_i$, $j\ne i$. Then $P(I+E_{ji})-P=E_{ii}$, again a difference of invertibles. This works over every field, including $\F_2$. Select $d^2$ independent invertibles from this spanning family over $k$ and lift their entries to $K$. Each lifted matrix is invertible. The square matrix of their flattened vectors is also invertible, because its residue is a basis matrix. The lifts therefore form the asserted basis.
\end{proof}

\begin{theorem}[Leading-layer orbit codebook]\label{thm:code}
For $M\in\Mat_d(K)\setminus\{0\}$, $d\geq2$, the maximum number of messages under the orbit-coding contract of Section~\ref{sec:contract} is $N_{\max}=dh$.
\end{theorem}
\begin{proof}
Write $M=\pi^aN$. The leading residue of each divided codeword $UN$ lies in
\[
 V=\{X\bar N:X\in\Mat_d(k)\},\qquad \dim_k V=dh.
\]
Each row of $X\bar N$ belongs to the $h$-dimensional row space of $\bar N$, and all such rows can be chosen independently, proving $\dim_kV=dh$. Put $v_i=\operatorname{vec}(U_iN)$ and let $D\in\GL_{d^2}(K)$ be a proposed decoder. The vectors $\bar v_i$ and $\bar D\bar v_i$ are nonzero. If a coordinate of $Dv_i$ has nonzero residue, it is a unit, and multiplication by $\pi^a\ne0$ cannot kill it. Thus
\[
\operatorname{supp}(\bar D\bar v_i)
\subseteq\operatorname{supp}(\pi^aDv_i).
\]
Exact decoding requires the supports on the right to be pairwise disjoint. Consequently the nonzero field vectors on the left have disjoint supports and are linearly independent. They lie in the $dh$-dimensional space $\bar D\operatorname{vec}(V)$, giving the upper bound. The reverse support inclusion is neither asserted nor needed.

By Lemma~\ref{lem:span}, images of invertible residue matrices span $V$. Choose $dh$ of these images that are independent and lift the encoders to $U_1,\ldots,U_{dh}\in\GL_d(K)$. The columns $\operatorname{vec}(U_iN)$ are independent modulo $J$. Extend the residue columns to a basis and lift the added columns arbitrarily. The resulting matrix $B$ over $K$ is invertible. Its inverse sends the actual codewords to $\pi^ae_i$, each with a unique possible coordinate outcome. This attains the bound.
\end{proof}

\begin{corollary}\label{cor:equiv}
With the complete-basis support and orbit-coding contracts matched, $N_{\max}>d$ if and only if $M$ is strongly contextual.
\end{corollary}
\begin{proof}
Both conditions are equivalent to $h\geq2$.
\end{proof}
The $d$-message comparison is the dimension baseline for a fresh $d$-coordinate carrier with complete basis readout. It is not an entropy or an asymptotic channel capacity. If encoders may discard the resource, or decoders are restricted, the stated optimization changes. Over a field $h=r$, giving the familiar linear-algebra orbit count $dr$; the two-mobit four-message protocol is already known \cite{Schumacher2010}.

A joint effect is a primitive covector
\[
e:K^d\otimes_K K^d\longrightarrow K,
\qquad E_{ij}=e(e_i\otimes e_j),
\]
where the first factor holds the input and the second Alice's share of the resource. A complete analysis is a basis of $d^2$ such covectors. Its analysis matrix has rows $\operatorname{vec}(E)^T$ and must be invertible. This is different from requiring an individual reshape $E$ to have matrix rank one: $E$ may be invertible and need not represent a decomposable local effect. For $M=\sum_{j,k}M_{jk}e_j\otimes e_k$, contraction gives
\[
(T_ex)_k=\sum_{i,j}E_{ij}x_iM_{jk},\qquad T_e=M^TE^T.
\]
The branch label is communicated to Bob, who applies a fixed correction for that label.

\begin{theorem}[Universal raw teleportation]\label{thm:teleport}
Let $K$ be a finite commutative local ring. For nonzero $M\in\Mat_d(K)$, $d\geq2$, universal exact raw-vector teleportation with a complete joint covector basis and unrestricted invertible corrections is possible if and only if $M$ is invertible.
\end{theorem}
\begin{proof}
Use the contraction convention above, so a branch with reshape $E$ has transfer $T=M^TE^T$. If a nonzero branch map had a nonzero kernel vector $z$, choose $x$ with $Tx\ne0$. Both $x$ and $x+z$ are nonzero and give the same possible branch output. A fixed correction cannot recover both. Thus each nonzero universally correct branch is injective. On the finite free module $K^d$, it is bijective and its inverse is linear; its invertibility forces $M$ to be invertible. A complete effect basis and nonzero $M$ exclude the possibility that every branch map is zero: the reshapes span all matrices, so this would force $M^TE_{ij}=0$ for all $i,j$ and hence $M=0$.

Conversely, choose the invertible-matrix basis of Lemma~\ref{lem:span} as the reshapes of the joint effects. It is a complete analysis basis, and every transfer $M^TE^T$ is invertible. Correct with its inverse. The identity $T^{-1}T=I$ remains an identity after tensoring with the identity on an untouched reference, proving the reference-preserving statement as well.
\end{proof}

The analyzer existence step is a residue-lifting corollary of the classical spanning fact. It requires both invertibility of each reshape and invertibility of the full analysis matrix. Its inclusion removes a dimension restriction from sufficiency when all the required bases and corrections are admitted; it does not assert that a physically restricted gate set can implement them.

\begin{corollary}[A nonprimitive separator]\label{cor:separator}
For a chain ring of length $s>1$ and $0<a<s$, the resource $\pi^a I_d$ is strongly contextual and has $d^2$ exact orbit messages, but cannot teleport universally under Theorem~\ref{thm:teleport}.
\end{corollary}
\begin{proof}
Its leading multiplicity is $d$. It also annihilates every vector multiplied by $\pi^{s-a}$, so is singular. Apply Theorems~\ref{thm:trichotomy}--\ref{thm:teleport}.
\end{proof}
The smallest example is $\varepsilon I_2$ over $\F_2[\varepsilon]/(\varepsilon^2)$. The role of nonprimitivity is exact:
\begin{corollary}[Primitive-resource boundary]\label{cor:primitive}
For a primitive $M\in\Mat_d(K)$ over a finite commutative chain ring, under the stated contracts,
\[
N_{\max}=d^2\quad\Longleftrightarrow\quad M\text{ invertible}
\quad\Longleftrightarrow\quad\text{universal exact raw teleportation}.
\]
\end{corollary}
\begin{proof}
Primitivity gives $a=0$. By Theorem~\ref{thm:code}, maximal coding is equivalent to $h=d$, so all Smith factors are units. Apply Theorem~\ref{thm:teleport}.
\end{proof}
Thus the maximal-coding/nonteleportation separation requires nonprimitive resources and does not occur in the generic state class of \cite{deBeaudrap2014}. This does not equate strong contextuality with teleportation for primitive resources: $2\leq h<d$ is still possible when $d>2$.

\section{Boundaries of the resource comparison}\label{sec:boundaries}
\subsection{Balanced tensors and regular binary expansion}
In $A=\A$, balanced scalar multiplication identifies $A\otimes_A A$ with $A$. Therefore
\[
 u^3\otimes_Au=0,\qquad u^3\otimes_{\F_2}u\ne0.
\]
The second assertion follows by applying field-linear functionals that are nonzero on each factor. Nonzero preparations thus need not remain nonzero under the balanced tensor. Primitive vectors do remain primitive, since a product of unit coordinates is a unit. Yet that restriction is not closed under the proposed conditional support update: applying the second standard effect to $\diag(1,u)$ gives the nonprimitive nonzero vector $(0,u)$.

The regular representation of $A$ is a faithful algebra map to $4\times4$ binary matrices. It is not a monoidal equivalence between the balanced and literal tensor spaces. Expanding a Smith factor $u^a$ gives binary rank $4-a$, whence
\[
 \rank_{\F_2}\lambda(M)=\sum_{i=1}^r(4-a_i).
\]
Unrestricted independent binary row and column operations classify a matrix only by that total rank. They can consequently erase distinctions detected by restricted $A$-linear operations.

For instance, $\diag(1,u^2)$ and $\diag(u,u)$ both have binary rank six but leading multiplicities one and two. To avoid relying on a nonprimitive comparison, the primitive matrices
\[
 \diag(1,1,u^3),\qquad \diag(1,u,u^2)
\]
both have binary rank nine, while their leading multiplicities are two and one. Their $A$-linear orbit codebooks have sizes six and three; both are singular, and hence neither universally teleports in the stated task.

\subsection{A boundary on the number of settings}
\begin{theorem}[Absence of strong contextuality with two settings]\label{thm:two}
A nonzero tensor with any finite number of parties over $\F_2$, two-dimensional local carriers, and at most two complete binary settings per party has a global assignment. The same conclusion holds over a finite commutative local ring with residue field $\F_2$.
\end{theorem}
\begin{proof}
Two distinct binary bases share a row $a$ and have other rows $b$ and $a+b$. If they coincide, use $b$ as the other row in both settings. Expand the tensor in the coordinates dual to $(a,b)$ at each site. Choose an inclusion-minimal set $S$ of sites using the $b$ coordinate whose coefficient is nonzero. Outside $S$ select $a$. Inside $S$ select the nonshared row of the chosen setting. Every joint contraction is the coefficient for $S$ plus coefficients of proper subsets, all of which vanish. Thus every selected joint event is supported.

For a local ring with radical $J$, choose the last nonzero radical layer containing all tensor coefficients: the tensor lies in $(J^a)^{2^n}$ but not in $(J^{a+1})^{2^n}$. Choose an $\F_2$-linear functional on $J^a/J^{a+1}$ for which the coefficient tensor is nonzero. Reduce each local basis modulo $J$ and apply the field argument. Every selected original contraction has nonzero image under this functional and is therefore nonzero.
\end{proof}
This does not exclude Hardy logical contextuality, which concerns failure to extend a particular event. Nor does it cover three settings or a different residue field. For two parties, the exclusion of a strongly contextual binary two-setting table is close to the established PR-box and modal-table boundaries \cite{AbramskyBrandenburger2011,Schumacher2012}; the proof above states explicitly what extends to arbitrary party number and radical layers.

\section{Finite verification and reproducibility}\label{sec:verification}
The proofs establish the general statements. The standalone supplement contains a separate checker developed for the September 2026 audit and rerun for this revision. It uses only the Python standard library. It is not the earlier historical runner, whose dependencies are absent from the original article package. The current evidence claims refer only to the files and commands in this supplement.

The support checker enumerates primitive rays, complete bases and exact contraction tables. Forbidden outcome pairs become binary clauses; implication-graph reachability decides global assignments and whether each supported event extends. Smith data are computed afterward for comparison, not used to decide support type. Identical support tables are cached. For the binary field these decisions are also checked by direct assignment enumeration.

\begin{center}\small
\begin{tabular}{p{.27\linewidth}p{.61\linewidth}}
\toprule
Check & Coverage and limitation\\\midrule
Direct support & All $22{,}170$ nonzero $2\times2$ resources over $\F_2,\F_3,\F_4$, $\F_2[u]/(u^2)$, $\F_2[u]/(u^3)$, $\F_3[u]/(u^2)$, and $\mathbb Z/4,\mathbb Z/8,\mathbb Z/9$; zero disagreements.\\[3pt]
Length-four support & Fourteen nonzero Smith representatives and 24 deterministic sampled resources over $A=\A$, each using all 192 local bases; not all $65{,}535$ nonzero resources.\\[3pt]
Hardy and analyzers & 36 ring/exponent Hardy cases; 28 analyzer configurations, checking each reshape and the complete analysis matrix, plus correction identities for one invertible resource per configuration.\\[3pt]
Coding & 52 explicit encoder/decoder certificates across seven rings, including all fourteen nonzero $A$ Smith representatives. These certify achievability; the general upper bound is analytic.\\[3pt]
Binary coding optimum & All 15 nonzero $2\times2$ binary resources, all $20{,}160$ invertible joint decoders per resource, and all disjoint-support subsets of each decoded orbit.\\[3pt]
Two-setting boundary & All 255 nonzero three-party binary tensors for a representative pair of distinct settings per party; the arbitrary-party and ring statements are analytic.\\\bottomrule
\end{tabular}
\end{center}

\begin{table}[ht]\centering\small
\begin{tabular}{rrrrrrll}
\toprule
$(a,b)$ & Count & $r$ & $h$ & Binary rank & Messages & Support & Teleport\\\midrule
$(0,0)$ & 24,576 & 2 & 2 & 8 & 4 & Strong & Yes\\
$(0,1)$ & 18,432 & 2 & 1 & 7 & 2 & Logical & No\\
$(0,2)$ & 9,216 & 2 & 1 & 6 & 2 & Logical & No\\
$(0,3)$ & 4,608 & 2 & 1 & 5 & 2 & Logical & No\\
$(0,4)$ & 4,608 & 1 & 1 & 4 & 2 & Local & No\\
$(1,1)$ & 1,536 & 2 & 2 & 6 & 4 & Strong & No\\
$(1,2)$ & 1,152 & 2 & 1 & 5 & 2 & Logical & No\\
$(1,3)$ & 576 & 2 & 1 & 4 & 2 & Logical & No\\
$(1,4)$ & 576 & 1 & 1 & 3 & 2 & Local & No\\
$(2,2)$ & 96 & 2 & 2 & 4 & 4 & Strong & No\\
$(2,3)$ & 72 & 2 & 1 & 3 & 2 & Logical & No\\
$(2,4)$ & 72 & 1 & 1 & 2 & 2 & Local & No\\
$(3,3)$ & 6 & 2 & 2 & 2 & 4 & Strong & No\\
$(3,4)$ & 9 & 1 & 1 & 1 & 2 & Local & No\\
$(4,4)$ & 1 & 0 & 0 & 0 & n/a & Excluded & n/a\\
\bottomrule\end{tabular}
\caption{All fifteen Smith classes over $A=\mathbb F_2[u]/(u^4)$ for $2\times2$ matrices. Exponent $4$ denotes zero. Counts are exhaustively enumerated; the remaining columns apply the analytic formulas. Logical means logical but not strong. Zero is excluded from the resource tasks.}\label{tab:smith}
\end{table}


Table~\ref{tab:smith} uses a separate exhaustive Smith inventory of all $65{,}536$ matrices over $A$, including zero. Its support labels are assigned by Theorem~\ref{thm:trichotomy}; they are not direct support decisions on the entire population. The nonzero counts are $5{,}265$ local, $34{,}056$ logical but not strong, and $26{,}214$ strong. The direct dual-number check yields $81$, $72$ and $102$ respectively for its 255 nonzero matrices.

A separately supplied process-semantics implementation is included unchanged for an additional, scoped cross-check: its Smith types and independently computed regular-representation binary ranks are checked on all $65{,}536$ matrices over $A$. Its inventory agrees with the audit checker. This comparison does not constitute another full support classification or a rerun of every historical process-semantics suite.

From the extracted revision directory, run \texttt{python supplement/run\_reproduction.py}. The wrapper refuses optimized Python execution, runs the full stated population, compares substantive outputs with the supplied reference results, regenerates Table~\ref{tab:smith}, and writes commands, interpreter/platform, exit statuses, runtimes and SHA-256 hashes. Runtime and platform fields are excluded from substantive-output comparisons. The Boolean coefficient-support identity is analytic; no historical LM-product or missing canonical-runner result is claimed as reproduced here. These finite checks corroborate the proofs and do not certify historical priority.

\section{Discussion}
The resource hierarchy is controlled by different questions about the same Smith form. At least two nonzero Smith factors permit a logical obstruction; leading multiplicity controls the absence of every global assignment and the size of an exact orbit code; invertibility controls universal raw-vector transfer. None of these statements identifies a physical capability without the relevant operation and readout assumptions.

There are two particularly useful limits. First, the nilpotent maximal-code example is a static nonprimitive resource, so it cannot by itself justify a closed process theory. For primitive resources, maximal orbit coding and universal raw teleportation coincide by Corollary~\ref{cor:primitive}. Second, the codebook/contextuality equivalence compares matched complete-basis contracts. Restricting gates, changing the scalar tensor, adding discard-and-reprepare, or weakening exactness requires a new theorem.

The next mathematical questions are a classification beyond chain rings, a process semantics closed under both preparation and conditioning, and resource comparisons for restricted measurement families. These remain open here. The contribution is the classification and the explicit comparison of resource criteria. The primary-source comparison is bounded; it does not certify historical firstness for the specific classification or orbit formula.

\bibliographystyle{plain}
\bibliography{references}
\end{document}
